% ==================================================================
%  ClientPreprocessing.tex  --  Section 3.1 "Client-side Preprocessing"
%  Loaded from main.tex.  Loads RoundCount.tex.
% ==================================================================

\subsection{Client-side Preprocessing}
\label{sec:client}

\subsubsection{Normalization}
The client first fixes one scale for the whole dataset. Let $g_{\min}$
and $g_{\max}$ be the smallest and the largest coordinate over all points
and all dimensions, and let
\[
  \sigma = g_{\max} - g_{\min} .
\]
The client maps every coordinate $x$ of every point to
\[
  x' = \frac{x - g_{\min}}{\sigma} \in [0,1],
\]
and the radius to
\[
  \varepsilon_n = \frac{\varepsilon}{\sigma} .
\]
Every dimension is scaled uniformly by the same $\sigma$, so every
pairwise distance shrinks by the same factor as the radius,
$\|\mathbf{x}' - \mathbf{y}'\| = \|\mathbf{x} - \mathbf{y}\| / \sigma$.
The values become smaller, but the clustering does not change:
$\|\mathbf{x} - \mathbf{y}\| \le \varepsilon$ if and only if
$\|\mathbf{x}' - \mathbf{y}'\| \le \varepsilon_n$. Every coordinate now
lies in $[0,1]$ and every squared distance in $[0,d]$, which gives the
bound on the comparison inputs required at the beginning of
Section~\ref{sec:method}.

\subsubsection{Ordering and the Offset Bound}
Deciding for every pair of points whether they are neighbors takes
$n(n-1)/2$ comparisons, and comparisons dominate the cost of the whole
protocol. Our goal is therefore to reduce the number of comparisons
itself.

Two choices achieve this. First, the server compares all pairs of slots
that are $k$ apart at once: for an offset $k$, a single evaluation
decides the $n$ pairs $(i, i+k)$ (Section~\ref{sec:neighbor}). Second,
the client orders the points so that neighbors lie in nearby slots, and
the server compares only the offsets at which neighbors can occur rather
than all of them.

The client obtains this order by PCA. A CKKS ciphertext holds the points
in a single row of slots, so the client arranges them along one
dimension: it projects every point onto the first principal component
$\mathbf{v}$ of the normalized data, $\|\mathbf{v}\| = 1$, sorts the
points by their projections $z_0 \le z_1 \le \dots \le z_{n-1}$, and
places the $i$-th point of this order in slot~$i$. The sorted vector keeps
neighbors close. The projection is $1$-Lipschitz by the Cauchy--Schwarz
inequality, so two points within distance $\varepsilon_n$ of each other
in $\mathbb{R}^d$ also lie within $\varepsilon_n$ of each other on the
sorted vector:
\[
  |z(\mathbf{x}) - z(\mathbf{y})|
  = |(\mathbf{x} - \mathbf{y})^{\top}\mathbf{v}|
  \le \|\mathbf{x} - \mathbf{y}\| \le \varepsilon_n ,
\]
where $z(\mathbf{x})$ denotes the projection of $\mathbf{x}$.

Consider the point in slot~$i$. Every neighbor in a later slot $j > i$
satisfies $z_j - z_i \le \varepsilon_n$ by this inequality, and
since the values are sorted, these slots form one run
$i+1, \dots, i+w_i$, where
\[
  w_i = \bigl| \{\, j > i : z_j - z_i \le \varepsilon_n \,\} \bigr| .
\]
Because the $z_i$ are sorted, the end of this run only moves forward as
$i$ grows, and all $w_i$ are found in a single pass in $O(n)$ time.

The server compares slots cyclically, modulo $n$
(Section~\ref{sec:notation}). Two slots that are more than $n/2$ apart
in one direction are therefore less than $n/2$ apart in the other, so
offsets up to $\lfloor n/2 \rfloor$, taken in both directions, already
reach every pair. The client records
\[
  k_{\max} = \min\Bigl( \bigl\lfloor n/2 \bigr\rfloor, \;
    \max_{i} w_i \Bigr).
\]

Comparing each slot only with the slots up to $k_{\max}$ away yields the
same clustering as comparing all $n(n-1)/2$ pairs. A pair that these
offsets leave out lies more than $\varepsilon_n$ apart on the sorted
vector, and hence, by the inequality above, more than $\varepsilon_n$
apart in $\mathbb{R}^d$. It is not a neighbor pair, so comparing it could
only produce a $0$ in an adjacency vector and would add nothing to any
neighbor count. Leaving it out changes the cost, not the result.

The bound $k_{\max}$ is the same for every point, the worst case over all
points, although a point in slot~$i$ needs only $w_i$ offsets. If the
server compared each point only with its own $w_i$ slots, the amount of
work spent on a point would follow the local density, and the server
could infer from it how the data are distributed, information that the
client never meant to reveal. A common bound also suits the slot-wise
evaluation, in which one operation acts on every point at once.

\input{\dirIII/RoundCount}

\subsubsection{Encryption}
The client encrypts one ciphertext per dimension: for
$j = 1, \dots, d$, the vector whose slot~$i$ is the $j$-th coordinate of
the point in slot~$i$, giving $\mathsf{ct}_1, \dots, \mathsf{ct}_d$. It
sends these to the server together with $\mathsf{pk}$, $\mathsf{evk}$,
and the public parameters $n$, $d$, $\varepsilon_n$, $\mathit{MinPts}$,
$k_{\max}$ and $R$, and keeps $\mathsf{sk}$ and the permutation that
produced the order. This is the only message the client sends.

% [2026-11-11] Summary removed to meet the 13-page limit; original text:
% \subsubsection{Summary of the Preprocessing}
% The client prepares its data with plaintext work only, in three steps,
% and then encrypts it.
% \begin{enumerate}
% \item \emph{Normalization} maps the data into $[0,1]^d$ so that every
%   comparison input on the server can be scaled into $[-1,1]$ by a public
%   constant. It outputs the normalized points and $\varepsilon_n$, in
%   $O(nd)$ time.
% \item \emph{Ordering} sorts the points along the first principal
%   component so that the server compares $k_{\max}$ offsets instead of all
%   pairs. It outputs the slot order and $k_{\max}$, in $O(nd^2 + d^3)$ time
%   for the principal component, $O(n \log n)$ for the sort and $O(n)$ for
%   $k_{\max}$.
% \item \emph{Round count} fixes how many rounds of label propagation the
%   server runs. It outputs $R$, in $O(n + T\,C \log C)$ time when the grid
%   estimate is used.
% \end{enumerate}
% \emph{Encryption} finally produces the $d$ ciphertexts
% $\mathsf{ct}_1, \dots, \mathsf{ct}_d$. For a fixed dimension $d$, the
% preprocessing therefore takes $O(n \log n + T\,C \log C)$ time, which is
% $O(n \log n + C^2 \log C)$ in the worst case; a client that sets $R$
% directly saves the grid term.
